\section{Introduction}

The proliferation of large language models (LLMs) has transformed the landscape of artificial intelligence, yet their development often relies on a legally and ethically precarious foundation. The vast majority of performant models are pretrained on massive web scrapes, indiscriminately mixing public-domain content with copyrighted materials such as books, news articles, and personal websites without explicit permission \citep{raffel2020exploring, gao2020pile}. This practice has led to a growing number of copyright infringement lawsuits,  creating significant legal uncertainty for both
academic researchers and commercial developers and threatening the future of the field. At the same
time, practitioners who wish to avoid this risk have few alternatives, as most high-performing
pretraining mixtures rely, at least in part, on opaque or non-permissive web scrapes.


Compounding this uncertainty is the prevailing assumption that state-of-the-art performance is inextricably linked to the sheer scale and diversity offered by these legally ambiguous web scrapes. The absence of a high-performance, large-scale pretraining dataset that actively mitigates these risks has forced a difficult choice between performance and compliance. In practice, the strongest open
baselines such as FineWeb-Edu~\citep{penedo2024} and DCLM~\citep{li2024datacomplm} still rely on mixed-license or unspecified web data,
whereas strictly permissive corpora tend to lag behind them on reasoning-heavy benchmarks.
This raises a critical question: Can
a powerful language model be trained on a dataset that provides a more legally robust foundation?

To this question, we answer "yes": We introduce \MixtureVitae{}, a \textbf{\DatasetSize{}}-billion-token, open-access pretraining dataset constructed to minimize copyright risk while explicitly demonstrating that a
reasoning- and instruction-dense, permissive-first mixture can substantially close the performance gap to leading
non-permissive corpora. The core of \MixtureVitae{}'s "permissive-first" data comprise (1) text with clear and permissive licenses (e.g., CC-BY-*, Apache 2.0), public-domain text, and copyright-exempt text such as US federal works (see Appendix~\ref{app:govt_works}) and (2) risk-mitigated text. Following Phi-4~\citep{abdin2024phi4technicalreport}, which shows that the addition of synthetic and web-rewrite data boosts performance, we address the scarcity of real, human-written reasoning and conversational dialogue in strictly permissive sources by significantly augmenting \MixtureVitae{} with targeted synthetic data, which is derived from permissive models and sources. We call this combination of expressly licensed and risk-mitigated methods the \textbf{"permissive-first"} approach. 


To validate our approach, we train models with \textbf{130M, 400M, 1.3B, and 1.7B parameters} on \MixtureVitae{} and compare their performance against several prominent open datasets. The results first confirm that \MixtureVitae{} \textbf{significantly outperforms all other permissively licensed baselines}, with the performance gap widening as the model scale increases. The more critical test, however, is against popular non-permissive datasets containing higher proportions of copyrighted or ambiguously-licensed material. In this setting, our models achieve competitive performance, and on math and code benchmarks, our 1.7B base model matches or exceeds a strong 1.7B instruction-tuned baseline despite being trained on a dramatically smaller token budget. 

In summary, our contributions are threefold:

\textbf{1. Permissive-first, risk-mitigated, and performant recipe for pretraining corpora.} We present \MixtureVitae{}, the first highly-performant, permissive-first, and risk-mitigated pretraining corpus that deliberately front-loads high-quality reasoning and instruction data to drive capability gains in small models. It is organized into auditable provenance tiers and constructed via a positive-inclusion pipeline, avoiding the need for retroactive filtering.

\textbf{2. We demonstrate that reliance on indiscriminately scraped, high-risk copyrighted data is not a prerequisite for training capable LLMs.} Leveraging the \texttt{open-sci-ref}~\citep{nezhurina2025opensciref001openreproduciblereference} protocol to ensure rigorous comparison across 130M–1.7B parameter scales, we demonstrate the value of front-loading instruction and reasoning data into pre-training. Our \DatasetSize{}B-token, permissive-first mixture closes the gap to mixed-license baselines while providing an auditable legal provenance. Furthermore, we show that our 1.7B base model, despite a limited 300B token budget, is comparable across multiple reasoning benchmarks to a strong 1.7B instruction‑tuned baseline—trained on roughly $36\times$ more tokens ($\approx$11T).

\textbf{3. Evaluation integrity and reusable artifacts}. We perform a large-scale 13-gram decontamination analysis across all benchmarks, showing that \MixtureVitae{}’s gains persist on decontaminated test sets and when removing shards responsible for most detected overlap, and we release the corpus, shard-level provenance metadata, and curation code to enable compliant, reproducible pretraining in future work. 

% \vspace{-0.2cm}

% \vspace{-0.4cm}

% \begin{itemize}
% \item
% \textbf{Permissive-first, risk-mitigated, and performant recipe for pretraining corpora.} We present \MixtureVitae{}, the first highly-performant, permissive-first, and risk-mitigated pretraining corpus that deliberately front-loads high-quality reasoning and instruction data to drive capability gains in small models. It is organized into auditable provenance tiers and constructed via a positive-inclusion pipeline, avoiding the need for retroactive filtering.
% \item
% \textbf{We demonstrate that reliance on indiscriminately scraped, high-risk copyrighted data is not a prerequisite for training capable LLMs.}
% Leveraging the \texttt{open-sci-ref}~\citep{nezhurina2025opensciref001openreproduciblereference} protocol to ensure rigorous comparison across 130M–1.7B parameter scales, we demonstrate the value of front-loading instruction and reasoning data into pre-training. Our \DatasetSize{}B-token, permissive-first mixture closes the gap to mixed-license baselines while providing an auditable legal provenance. Furthermore, we show that our 1.7B base model, despite a limited 300B token budget, is comparable across multiple reasoning benchmarks to a strong 1.7B instruction‑tuned baseline—trained on roughly $36\times$ more tokens ($\approx$11T).

% \item
% \textbf{Evaluation integrity and reusable artifacts}. We perform a large-scale 13-gram decontamination analysis across all benchmarks, showing that \MixtureVitae{}’s gains persist on decontaminated test sets and when removing shards responsible for most detected overlap, and we release the corpus, shard-level provenance metadata, and curation code to enable compliant, reproducible pretraining in future work.

% \end{itemize}
